\section{\ENG{Player Sandbox} και \ENG{Security Boundaries}}
\label{sec:player_sandbox_security}

Το περιεχόμενο εκτελείται σε περιβάλλον \ENG{browser}, άρα η ασφάλεια δεν μπορεί να στηριχθεί μόνο στην αυθεντικοποίηση των \ENG{backend} υπηρεσιών. Η υλοποίηση αντιμετωπίζει το \ENG{SCO} ως δυνητικά μη έμπιστο περιεχόμενο και το απομονώνει μέσα σε \ENG{iframe}. Το \ENG{launch page} λειτουργεί ως \ENG{host} περιβάλλον, εκθέτοντας μόνο τον ελάχιστο \ENG{API bridge} που απαιτείται για \ENG{SCORM} επικοινωνία και διατηρώντας τη δική του λογική \ENG{commit/terminate} εκτός του χώρου εκτέλεσης του ίδιου του περιεχομένου.

Το πρώτο όριο ασφάλειας αφορά την πρόσβαση στα αρχεία περιεχομένου. Η κανονικοποίηση της σχετικής διαδρομής, η παραγωγή απόλυτης διαδρομής εντός του \ENG{launch} καταλόγου και η απόρριψη κάθε αιτήματος που επιχειρεί να βγει έξω από αυτόν μέσω \ENG{path traversal} διασφαλίζουν ότι ο \ENG{route} \texttt{\ENG{/content/:launchId/*}} δεν μπορεί να χρησιμοποιηθεί για αυθαίρετη ανάγνωση αρχείων του συστήματος αρχείων.

Το δεύτερο όριο αφορά τη γονική εφαρμογή \ENG{LMS}. Όταν ο \ENG{player} ολοκληρώνει μια συνεδρία, στέλνει μήνυμα \texttt{\ENG{SCORM\_PLAYER\_TERMINATED}} προς το γονικό παράθυρο. Η σελίδα \texttt{\ENG{learner/launch.blade.php}} δεν αποδέχεται οποιοδήποτε \ENG{postMessage}, αλλά ελέγχει πρώτα αν το \ENG{origin} του συμβάντος ταιριάζει με το \ENG{origin} του \ENG{player iframe}.

Το τρίτο όριο είναι λειτουργικό και αφορά τη σχέση ανάμεσα στον \ENG{player} και στο \ENG{engine}. Το \ENG{SCO} δεν αποκτά άμεση πρόσβαση στα \ENG{launch tokens} πέρα από το απολύτως αναγκαίο \ENG{bridge} που έχει συνθέσει η σελίδα του \ENG{player}. Οι πραγματικές \ENG{HTTP} κλήσεις που φθάνουν στο \ENG{engine} κατασκευάζονται από τον \ENG{host} κώδικα του \ENG{player}, όχι από το ίδιο το \ENG{SCO}. Η επιλογή αυτή δεν μηδενίζει τους κινδύνους ενός κακόβουλου περιεχομένου, αλλά περιορίζει τα εκτεθειμένα \ENG{endpoints} και κρατά τον αυστηρότερο έλεγχο στον \ENG{server-side} μηχανισμό του \ENG{player}.

Το Σχήμα~\ref{fig:player_sandbox_security} απεικονίζει τα ανωτέρω όρια. Η σημασία του δεν είναι απλώς διαγραμματική. Δείχνει ότι η \ENG{client-side} αρχιτεκτονική του συστήματος δεν είναι ταυτόσημη με εκείνη μιας κοινής \ENG{file server} προσέγγισης, αλλά στηρίζεται σε ένα ελεγχόμενο \ENG{host layer} που απορροφά το μεγαλύτερο μέρος της διαμεσολάβησης ανάμεσα στο περιεχόμενο και στην πλατφόρμα.

\begin{figure}[H]
  \centering
\begin{chapterfigurebox}
  \begin{tikzpicture}[
      font=\footnotesize,
      node distance=1.9cm and 2.5cm,
      box/.style={draw, rounded corners, fill=gray!5, align=center, text width=3.1cm, minimum height=1.05cm},
      emph/.style={draw, rounded corners, fill=gray!12, align=center, text width=3.5cm, minimum height=1.05cm, font=\footnotesize\bfseries},
      arrow/.style={-Latex, thick},
      boundary/.style={draw, dashed, rounded corners, inner sep=11pt}
    ]
    \node[emph] (host) {\ENG{Player host page}};
    \node[box, right=of host] (sco) {\ENG{SCO iframe}};
    \node[box, left=of host] (parent) {\ENG{LMS parent}\\window};
    \node[emph, below=2.0cm of host] (engine) {\ENG{SCORM Engine}};

    \draw[arrow] (sco) -- node[above, align=center]{\ENG{SCORM API}\\\ENG{bridge}} (host);
    \draw[arrow] (host) -- node[left, align=center]{\ENG{commit /}\\\ENG{terminate}} (engine);
    \draw[arrow] (host) -- node[above]{\ENG{postMessage}} (parent);

    \node[boundary, fit=(host)(sco)] (browserboundary) {};
    \node[font=\footnotesize, align=center, text width=5.1cm, above=0.14cm of browserboundary] {\ENG{Browser execution boundary}};

    \draw[dashed] ($(host.east)!0.5!(sco.west)$) ++(0,0.9) -- ++(0,-2.0);
    \node[font=\footnotesize, align=center] at ($(host.east)!0.5!(sco.west)+(0,1.25)$) {\ENG{iframe isolation}};
  \end{tikzpicture}
\end{chapterfigurebox}
  \caption{Κύρια όρια ασφάλειας του \ENG{player}: απομόνωση \ENG{iframe}, \ENG{host-mediated} πρόσβαση στο \ENG{engine} και ελεγχόμενο \ENG{postMessage} προς το \ENG{LMS}.}
  \label{fig:player_sandbox_security}
\end{figure}

